\section{Beyond Traditional Data: Text, Networks, Images, and Treatment over Time}
\label{section:nontraditional}

As exciting as neural networks are for heterogeneous treatment effect estimation from quantitative data, a great promise of deep causal estimation is inference when treatments, confounders, and mediators are encoded in high-dimensional data (e.g., text, images, social networks, speech, and video) or are time-varying. This is a strong advantage of neural networks over other machine learning approaches, which do not generalize competitively to non-quantitative data.  In these scenarios, multi-task objectives and tailored architectures can be used to learn representations that are simultaneously rich, capture information about causal quantities, and disentangle their relationships. Moreover, the inherent flexibility of neural networks means that, in many cases, the TARNet-style models presented above can serve as the foundations to inference on text and graphs with some architectural modifications, additional losses, and new identification assumptions. 

This literature is rapidly evolving, so readers should treat this section of the primer as fundamentally prospective. To maintain accessibility, our primary goal here is to introduce readers to hypothetical scenarios where they might perform causal inference on text, network, or image data. Second, we selectively review contemporary, theoretically-motivated literature on deep causal estimation in these settings. The identification assumptions for different data types differ substantially, so we generally leave those to the interested reader. Finally, we briefly discuss approaches for dealing with time-varying confounding. We also take this section as an opportunity to introduce the Transformer or Graph Neural Network, an architecture now used in most contemporary deep learning models to learn from complex data (Box \ref{Box:gnn}).
 
\subsection{Causal Inference from Text}
\label{subsection:text}

In recent years, an interdisciplinary community across both social science and computer science has coalesced around causal inference from text (see \citet{keith-etal-2020-text} and \citet{feder2021causal} for exhaustive reviews). Broadly speaking, texts may capture information about any causal quantity (treatments, outcomes, confounders, mediators) we might be interested in. For example, in an exit-polling experiment, analysts might want to measure toxicity ($Y$) in text responses to political prompts. In an observational study of e-mail response times ($Y$), analysts might want to measure the effects of the tone of the email ($T$). In this scenario, the analyst might also want to control for confounders like subject matter ($X$). Each of these scenarios presents distinct identification challenges \citep{feder2021causal}. But in all cases, we can use low-dimensional representations of the high dimensional text to extract, quantify, and disentangle relationships between nuanced qualities like tone and subject matter. %Research within the social science community has focused on representations created by Bayesian topic models for these problems, while computer scientists have predominantly explored deep learning solutions.

The ability of neural networks to automatically extract features makes them particularly suited for the last scenario when both treatment information and confounding covariates are encoded in text. In many cases, we may not have explicitly identified, quantified, or labeled all of the confounders in text (e.g., subject matter and tone of emails), but we would still like to control for them. \citet{pryzant2020causal},\citet{Veitch2019UsingTE}, and \citet{gui2022causal} address this problem by prepending Transformer-layers (Box \ref{Box:gnn}) for reading text to the beginning of TARNet or Dragonnet. \citet{Veitch2019UsingTE} demonstrate the viability of this approach on a Science of Science question testing the causal effect of equations on getting papers accepted to computer science conferences. \citet{pryzant2020causal,gui2022causal} explore the more complicated scenario not in which the treatment is explicitly known (e.g., equations in papers, gender of authors), but is instead externally perceived upon reading (e.g., politeness/rudeness of an email or toxicity of a social media post). In these models, an additional loss function is also added for learning text representations concurrently with the causal inference losses discussed above. 

 \begin{TCBBox}[label=Box:gnn]{Graph Neural Networks and Transformers}\
 \textbf{Graph neural networks} (GNNs) are the current state-of-the-art approach for creating representations for nodes in graphs. Compared to previous approaches that relied on ``shallow" embeddings based only on a node's local context (e.g., random walks to nearby nodes), GNNs are attractive because their node representations are aggregated from the structural position and covariates of all nodes $n$ degrees away from the target node, where $n$ is the number of graph neural network layers.


The most intuitive understanding of how graph neural networks work is as a message passing system \citep{gilmer2017neural}. We use one of the first GNN papers, the Graph Convolutional Network as an example \citep{kipf2016semi}. In this interpretation, each node has a message that it passes to it's neighbors through a graph convolution operation. In the first layer of a GNN this message would consist of the node's covariates/features. In consecutive layers of the network, these messages are actually representations of the node produced by the previous layer. During graph convolution, each node multiplies incoming messages by it's own set of weights and combines these weighted inputs using an aggregation function (e.g., summation). By the $n$-th GNN layer, these messages will contain structure and covariate information from all nodes $n$ degrees away. For interested readers, there is also a spectral interpretation of this process. Typically GNNs are trained to produce representations of graphs by predicting the probability that two nodes are linked in the network, and then used for something else. One variant of the GNN uses an ``attention" mechanism to vary the extent that nodes value messages from different neighbors (the graph attention network or GAT) \citep{velivckovic2017graph}. 

As of 2023, \textbf{Transformers} are the hegemonic architecture used in natural language processing. After their introduction in 2017, these models improved performance on many high-profile NLP tasks across the board. Several enterprise-scale transformers have been featured in the media for their impressive performance in text generation and question answering (e.g. OpenAI's GPT-3 and ChatGPT, Google's Bard). Smaller models in broad use are based on the BERT architecture \citep{devlin2018bert}.

The connection between GNNs, and specifically GATs, is the focus on attention mechanisms. From this perspective, words in sentences are akin to nodes in networks, with their relative positions to each other being analogous to their structural positions in the graph. Transformers improved on previous sequential approaches to text analysis (i.e. recurrent neural networks) by having each word (or representation of a word) receive messages from not just adjacent words, but all words heterogeneously. Attention mechanisms throughout the architecture allow each layer of a transformer to attend to words or aggregated representation mechanisms heterogeneously. Architectures such as BERT or GPT stack transformer layers to create models with hundreds of millions of parameters. These models are expensive to train, both computationally and with respect to data, so they are often pretrained on large datasets and then ``fine-tuned" (lightly re-trained) with smaller datasets for specific tasks or to align with certain goals.  

\end{TCBBox}

\subsection{Causal Inference from Networks}
\label{subsection:networks}

A smaller literature has leveraged relational data for causal inference in two distinct scenarios. In the first traditional selection on observable settings, we wish to control for information about unobserved confounding inferable from homophilous ties. For example, age or gender might be unmeasured in our data, but we might expect people to develop friendship ties with those of the same gender identity or age cohort.

This scenario suggests estimation strategies similar to those when confounders are encoded in text. Much like Transformer layers can be prepended to TARNet-style estimators to learn from text, graph neural networks (an analog of the Transformer) can be preprended to learn from graphs. \citet{guo2019counterfactual} provides a first pass at this problem by adding  GNN layers to CFRNet \cite{Shalit2017}. \citet{Veitch2019UsingET} instead adapt Dragonnet in a semi-parametric framework to allow for consistent estimates of the treatment and outcome, assuming the network representation encodes significant information about confounders.

The second, more challenging scenario is estimating the causal effect of social influence on outcomes from observational data. For example, \citet{Cristali2022UsingEF} introduce the problem of measuring the effects of vaccination ($T$) on peer vaccination choice ($Y$). This is a hard problem because a) SUTVA is a fundamental assumption of all causal inference frameworks and b) it is hard to disentangle whether changes in the outcome result from the treatment via peer effects (e.g, person A pressuring person B to vaccinate), or from homophily (e.g., person A and person B having similar political leanings). In other words, contagion and homophily are generically confounded \citep{shalizithomas2011}. \citet{mcfowland2021} are the first to tackle this problem by making strong parametric assumptions about the generation of network ties and the outcome model. \citet{Cristali2022UsingEF} instead propose an approach using neural network-learned representations of the graph.     

\subsection{Causal Inference from Images}
\label{subsection:images}

While ideas from causal inference have been leveraged extensively to improve image classification, to our knowledge there are no papers that explore causal inference where treatments, confounders, mediators, or predictors are encoded in images.\footnote{\citet{jesson2021quantifying} introduce a simulation where the MNIST digit dataset serves as covariates $X$ as toy example of high-dimensional confounding, but not a possible application.} That being said, some scenarios proposed for causal text analysis should apply here as well. For example, consider the conjoint experiment by \citet{todorov2005inferences} where both the treatment (e.g. incumbency of a politician) and potential latent confounders (e.g., party, age, gender, race) are encoded in an image. In this setting, a TARNet-like model adapted to learn and condition on image representations could improve treatment effect estimation by controlling for confounders such as the politician's age. Causal inference on images is an area ripe for exploration, and we hope to see more work here in the future.

\subsection{Causal Inference from Time-varying Data}
\label{subsection:timevarying}

One natural extension of deep causal estimation is to scenarios where treatments are administered over time and confounding may be time-varying. While “g-methods” developed by Robins et al. for estimating effects with time-varying treatments and confounding have existed for decades, the statistical assumptions encoded in these models are quite strong \citep{robins1994correcting,robins2000marginal,robins2009longitudinal}. Due to their reliance on generalized linear models to define the “structural” component, they assume that the outcome is a linear function of all covariates and treatment. Second, for identification, they make strong assumptions about which previous timesteps confound the current one. Third, they require different coefficients to be estimated at each time steps. Transformers (Box \ref{Box:gnn}) and recurrent neural networks, a simpler model for sequential data (Appendix A.\ref{appendix:rnn}), should be able to capture long-term dependencies and non-linearities in ways that marginal structural models and g-computation cannot.

Several papers have begun to explore these possibilities in the context of personalized medicine. \citet{lim2018forecasting} build a marginal structural model using a recurrent neural network, and \citet{bica2020estimating} extend this framework with an additional loss to more explicitly deal with time varying confounding by forcing the model to “unlearn” information about the previous time steps. \citet{melnychuk2022} go one step further by adapting \citet{bica2020estimating}’s approach with a transformer. Inspired by longitudinal targeted maximum likelihood, \citet{frauen2022estimating} add a semi-parametric targeting layer to their RNN to create a g-computation algorithm that is doubly robust and asymptotically efficient. \citet{lietal2021} instead propose an RNN framework for g-computation that allows for dynamic treatment regimes. All of these papers use simulations of tumor growth dynamics, naturalistic simulations based on vital signs from intensive care unit visits, or factual datasets exploring treatment response to physical therapy for back pain.